\tableofcontents
\vspace{1em}
\hrule
\vspace{1em}

\section{Some Basics}

The tasks are synthetic retrieval (RULER): find a random number planted somewhere in 8K–128K tokens of text, with distractor numbers nearby. There are 20 prompts per cell.
\begin{itemize}
    \item Synthetic Retrieval: This means the test is artificially constructed rather than being a natural real-world task (like summarizing a real book). A piece of data is intentionally hidden inside a massive wall of text just to see if the AI can find it. 
    \item RULER: This is the specific name of the benchmark framework (developed by researchers to evaluate long-context LLMs). It stands for Retrieval, Understanding, and Long-context Evaluation Robustness.
    \item 8K-128K tokens: This indicates the length of the text context that the model has to process. A token is roughly a word or part of a word, so 8K tokens is about 8,000 words, and 128K tokens is about 128,000 words. The model has to search through this large amount of text to find the hidden number.
    \item Distractor numbers: These are other numbers placed in the text that are not the target number. They are meant to confuse the model and make the task harder. The model has to correctly identify the planted number and ignore the distractors.
    \item 20 prompts per cell: The Matrix (Cells): When testing LLMs, researchers usually create a grid or matrix. For example, the horizontal axis might be the context length (8K, 16K, 32K, 64K, 128K) and the vertical axis might be the depth of where the needle is hidden (0\% depth = top of the document, 50\% = middle, 100\% = bottom). Each intersection on this grid is a "cell." This means that for each configuration or "cell" of the experiment (which could be a specific model, context length, or other parameter), there are 20 different test cases (prompts) that the model is evaluated on. Each prompt is a separate instance of the retrieval task. 
\end{itemize}

Some metrics
\begin{itemize}
    \item valid mean task score
\end{itemize}

10 promots -> 100prompts or 1000 prompts? 
not  head-output error but task score.  The task score is the final evaluation metric that measures how well the model performed on the retrieval task. It is calculated based on whether the model correctly identified the planted number in each prompt. The valid mean task score is the average score across all valid prompts (those that were not filtered out due to errors or other issues). This metric gives an overall sense of the model's performance on the retrieval task across different configurations and contexts.


\section{Scope of the common rate-allocation view}
\label{app:scope}

This appendix expands the derivation and measurement contract used in the main
paper.  The common action space is deliberately narrow: it allocates
\emph{key-token} rates within an attention head, leaves values exact, and measures
attention-output error.  Prior work already combines quantization and eviction
in deployable systems.  Our claim is instead that a common local surrogate makes
it possible to ask a different empirical question: when does the discrete
interior improve over uniform quantization and eviction under matched
information and key-bit budgets?

Throughout, \emph{zero rate} means a zero-bit action in the separable planning
surrogate.  It does not mean that set eviction is exactly identical to
independent additive logit noise.  Section~\ref{app:derivation} makes this
distinction explicit.

\section{Output distortion and the zero-rate approximation}
\label{app:derivation}

\subsection{Attention and finite-bit key perturbations}

For a query $q\in\mathbb R^d$, keys $k_i\in\mathbb R^d$, and values
$v_i\in\mathbb R^{d_v}$, define
\begin{align}
 s_i &= q^\top k_i/\sqrt d, &
 a_i &= \frac{\exp(s_i)}{\sum_j\exp(s_j)}, &
 o &= \sum_i a_i v_i .
\end{align}
Quantizing key $i$ changes its logit by $\delta_i$.  Because
\[
 \frac{\partial o}{\partial s_i}=a_i(v_i-o),
\]
the first-order output perturbation is
\begin{equation}
 \Delta o \approx \sum_i a_i(v_i-o)\delta_i .
 \label{eq:app-linear}
\end{equation}
Consequently,
\begin{align}
 \mathbb E\|\Delta o\|^2
 &\approx
 \sum_i a_i^2\|v_i-o\|^2\operatorname{Var}(\delta_i)
 \nonumber\\
 &\quad+
 \sum_{i\ne j}a_i a_j
 (v_i-o)^\top(v_j-o)\operatorname{Cov}(\delta_i,\delta_j).
 \label{eq:app-covariance}
\end{align}
Equation~\eqref{eq:surrogate} follows only when the residual perturbations are
zero mean and the cross-token covariance contribution vanishes.  Softmax
invariance removes a \emph{common} shift,
$\operatorname{softmax}(s+c\mathbf 1)=\operatorname{softmax}(s)$; it does not
remove arbitrary token-dependent bias.

\subsection{Exact set eviction}

Let $E$ be an evicted token set and $A_E=\sum_{i\in E}a_i$.  Renormalizing the
remaining attention weights yields
\begin{align}
 o_{\setminus E}
 &=\frac{o-\sum_{i\in E}a_i v_i}{1-A_E},\\
 o_{\setminus E}-o
 &=-\frac{\sum_{i\in E}a_i(v_i-o)}{1-A_E},\\
 \|o_{\setminus E}-o\|^2
 &=\frac{\left\|\sum_{i\in E}a_i(v_i-o)\right\|^2}
 {(1-A_E)^2}.
 \label{eq:app-exact-eviction}
\end{align}
The denominator and cross-token inner products prevent an exact token-separable
representation.  For one evicted token $j$, the exact error magnitude is
\[
 \frac{a_j}{1-a_j}\|v_j-o\|,
\]
which connects directly to the CAOTE singleton score.  ReST-KV further models
attention redistribution and output reconstruction; our local planner does not
subsume those effects.

\subsection{The planning approximation}

When the evicted mass is small and cross terms are neglected,
Eq.~\eqref{eq:app-exact-eviction} motivates
\[
 \widetilde D_0(E)=\sum_{i\in E} w_i,
 \qquad
 w_i=a_i^2\|v_i-o\|^2.
\]
We choose units in which the absolute zero-rate distortion is $d_0=1$ and write
the finite-rate costs as $d_b=\operatorname{Var}(\delta_b)$ in the same logit
units.  Thus
\[
 \widetilde D(\mathbf b)=\sum_i w_i d_{b_i},
 \qquad b_i\in\mathcal B,\quad 0\in\mathcal B.
\]
This convention keeps relative quantizer factors separate from the absolute zero-rate
cost.  It is a local approximation used to choose an allocation; every
reported response is obtained by recomputing attention with the actual
quantized/evicted cache.

\section{Discrete rate allocation and regime coordinates}
\label{app:allocation}

\subsection{The implemented optimizer}

At average key budget $B$, the discrete problem is
\begin{equation}
 \min_{b_i\in\mathcal B}\sum_i w_i d_{b_i}
 \quad\text{subject to}\quad
 \sum_i b_i\le BL.
 \label{eq:app-discrete}
\end{equation}
Introducing multiplier $\lambda\ge0$ separates the token decisions:
\begin{equation}
 b_i^\star(\lambda)
 =\arg\min_{b\in\mathcal B}\{w_i d_b+\lambda b\}.
 \label{eq:app-lagrange}
\end{equation}
We search $\lambda$ to satisfy the budget and use the discrete tiers in
Eq.~\eqref{eq:app-lagrange}; the continuous expression is explanatory, not the
experimental optimizer.

For the relaxation $d_b=\kappa 4^{-b}$,
\[
 b_i^\star
 =\frac12\log_2w_i+c
 =\log_2a_i+\log_2\|v_i-o\|+c.
\]
The measured $0.035$-bit result is the maximum observed change in an
\emph{aggregate ladder-width statistic} after adding the value term.  It is not
a uniform per-token bound and does not establish that attention-only allocation
is optimal.  A direct value-aware versus attention-only allocation ablation is
therefore retained as an explicit evidence anchor.

\subsection{Tier extinction}

If $d_b\ge d_0$, tier $b$ consumes positive bits without improving the local
surrogate relative to zero rate and is dominated.  This is a sufficient
extinction condition.  The converse is not sufficient: a tier with $d_b<d_0$
can still lie above the lower convex envelope of the other rate--distortion
points.  For the 2-bit tier we report
\[
 D_2=\frac1H\sum_{h=1}^{H}\mathbb I[d_{h,2}\ge d_0].
\]
$D_2$ is measured at a particular model and context.  Although $L$ is absent
from its formula, the head cost distribution changes with $L$, so $D_2$ is not
sequence-length invariant and a new deployment context requires calibration.

\subsection{Why the regime map needs two coordinates}

The sharp side and diffuse side arise for different reasons:
\begin{enumerate}
 \item On the sharp side, large absolute finite-bit noise removes low-bit tiers;
 $D_2$ directly probes this boundary.
 \item On the diffuse side, unequal allocation has little leverage when token
 importances are nearly equal; an importance-dispersion statistic is needed.
\end{enumerate}
For attention-only importance,
$\operatorname{std}_i(\log_2 a_i)=\tau/\ln2$, because softmax subtracts the same
log normalizer from every logit.  This identity alone neither proves that
$\tau$ is a complete regime variable nor establishes a uniform-quantization
optimum.  The paper therefore treats the mixed-to-eviction side as the
well-supported boundary and leaves a prospective comparison of dispersion
predictors as R10.

\section{Measurement protocol and accounting}
\label{app:protocol}

\subsection{Grid and independent units}

The symmetric R3 grid contains 16 model--context cells from six models and
29,312 query-head--configuration aggregates.  The R6 boundary extension contains
25 cells spanning 4k--256k and 44,416 such aggregates.  Across the six models,
10,240 query-head positions and 2,016 KV-head positions are counted once;
contexts reuse those positions.  Model/cell-level analyses, rather than raw
heads, are therefore the independent statistical units.

PG-19 passages supply natural context.  Each cell is admitted only after the
uncompressed model passes the needle-retrieval validity gate.  The allocation
planner uses its local surrogate, while every reported error comes from actual
key quantization or eviction followed by exact attention-output recomputation.

\subsection{Matched information and budget}

The primary comparison uses $B=3$ key bits per context token.  Uniform
quantization stores every key at 3 bits.  The eviction endpoint retains the
number of 8-bit keys allowed by the same key-bit budget.  In the symmetric
contract, the mixed allocator and accumulated-attention eviction both use lagged
importance information; an oracle/current-query result is labeled only as an
upper bound.  The five-corner robustness result changes the reference set and is
reported separately.

This accounting omits value-cache bytes, quantization scales, indices, metadata,
protected windows, alignment, and packing.  Moreover, \emph{8-bit} is the maximum
experimental key tier, not full floating-point precision.  The experiments
therefore establish matched key-bit output error, not matched total memory or
throughput.

\subsection{Responses}

For each head,
\[
 g_h=
 \frac{\min(D_h^{\rm uniform},D_h^{\rm evict})}
 {D_h^{\rm mixed}},
 \qquad
 G_+=\exp\!\left(\frac1H\sum_h\log\max(g_h,1)\right).
\]
$G_+$ is a hindsight portfolio statistic: it clips losses after observing all
three errors.  It is not the performance of a deployed router.  The primary
cell-level response is continuous $G_+$; the fraction with $g_h\ge2$ is a
secondary thresholded band statistic.  Every band result names its corner
and information contracts.

\section{Additional robustness and deployment evidence}
\label{app:robustness}

The completed measurements support the following bounded conclusions:
\begin{itemize}
 \item In all 16 R3 cells, the symmetric lagged-versus-lagged band exceeds the
 all-oracle band by 1.2--12.5 points, but it is lower than the asymmetric
 current-query-versus-lagged estimate.  Five cells are STOP and five NARROW.
 \item Across the 25 R6 cells, $D_2$ has Spearman correlation $-0.978$ with both
 $G_+$ and the band; the model-partial correlation is $-0.915$.  Similar $D_2$
 values can nevertheless differ by eight band points across architectures.
 \item Prompt-block standard deviation is 0.4--3.0 band points.  Replacing the
 two headline corners by the five-corner minimum lowers the band by 5.1--8.6
 points and can change a categorical verdict.
 \item In six R5 cells, the calibration-time family choice has p90 regret 1.00
 through 4,096 generated tokens, whereas a frozen token allocation costs
 $2.52$--$4.30\times$ a one-step-old allocation.
 \item Enforcing one allocation per shared KV head gives marginal in-band costs
 $1.000$, $1.136$, and $1.316$ at GQA repetition factors 1, 4, and 8.  A 4-bit
 current-query cascade closes 82--91\% (median 89\%) of the lagged-score penalty.
\end{itemize}

One predictor study remains intentionally unfilled.  R10 must compare $D_2$
prospectively against RateQuant-style heterogeneity, entropy/Gini, and
context-only predictors under a fixed leave-model-out protocol.  The end-task
audit is reported in Appendix~\ref{app:endtask}.  Beyond it, two natural-task
(LongBench-v2 forced-choice) development blocks expose only $3/52$ and $2/52$
candidate headroom over their best fixed policies; sequential gates suppress
both proxy outcomes and leave confirmation unopened.

\section{Claim boundary}
\label{app:claims}

\begin{table}[h]
\centering
\small
\begin{tabular}{@{}p{0.29\textwidth}p{0.64\textwidth}@{}}
\toprule
\textbf{Supported} & \textbf{Not yet supported} \\
\midrule
Zero rate is useful inside a stated separable local surrogate
& Exact equivalence between independent quantization noise and set eviction \\
An absolute tier cost can make a finite-bit tier dominated
& A universal one-dimensional phase transition or threshold \\
$D_2$ strongly orders the measured mixed-to-eviction boundary
& A context-free predictor that transfers without calibration \\
The output-error-optimal family label is temporally stable in six measured cells while token allocations become stale quickly
& A robust end-task router, total-KV memory savings, or realized throughput \\
The output-error-optimal allocator is the strongest token-selective policy on the measured retrieval cells
& That it beats dense low-bit quantization in general (it loses on Llama-3.1-8B) \\
Output distortion orders budgets and contexts by difficulty
& That output distortion can choose between method families at a fixed budget \\
Token-granular allocation truncates multi-token answers
& That positional smoothing repairs it (pre-registered test pending) \\
\bottomrule
\end{tabular}
\caption{The empirical and methodological claim boundary of the current draft.}
\label{tab:app-claims}
\end{table}

\section{End-task audit: protocol, baselines, and full results}
\label{app:endtask}

\paragraph{Protocol.}
Each prompt's context is prefilled at full precision and scored; the cache is
then compressed once, and the question is prefilled \emph{through} the
compressed cache before greedy decoding.  This is the question-agnostic setting
of a prefix cache or multi-turn session.  When the question is inside the
observation window instead, SnapKV scores 1.00 at every budget, because the
window vote points at the needle.  Evaluation uses prompts 100--119; the
\sieve{} router is calibrated on the disjoint prompts 0--9.  A block is valid
when the uncompressed model scores at least 0.9; this excludes Qwen3-8B
multi-value retrieval (0.24 at 8k, 0.04 at 32k).  Scores are RULER string
matches; multi-value and variable tracking give fractional credit.

\paragraph{Budget accounting.}
Every method spends $B$ key code bits per context-key element, verified at run
time from the applied bit maps.  Side information differs by method and is
reported in Table~\ref{tab:app-side} rather than folded into $B$.  Eviction and
\sieve{} keep tokens at 8 bits with the same TurboQuant quantizer, so their
side information scales with the kept fraction.

\begin{table}[h]
\centering
\small
\setlength{\tabcolsep}{4pt}
\caption{Baselines as implemented, and side information per key element ($d=128$).}
\label{tab:app-side}
\begin{tabular}{@{}p{0.2\linewidth}p{0.52\linewidth}r@{}}
\toprule
method & what is stored, and how it is chosen & side bits \\
\midrule
TurboQuant-MSE & every token at $B$ bits; random rotation, Lloyd--Max, per-token norm & 0.125 \\
KIVI ($g=32$ / $128$) & every token at $B$ bits; per-channel asymmetric integer, min/max per $g$ tokens, post-RoPE & 1.0 / 0.25 \\
KVQuant & every token at $B$ bits; pre-RoPE per-channel non-uniform codebook; 1\% outliers and token 0 exact & $\approx0.32$ \\
H2O & $\lfloor BC/8\rfloor$ tokens at 8 bits; attention summed over all prefill queries & $\approx0.05$ \\
SnapKV & same keep count; 32-query observation window, max-pool 7 & $\approx0.05$ \\
DropKV, Ada-KV, OBCache-K, LaProx & same keep count; published scores; Ada-KV head allocation for Ada-KV and OBCache-K; LaProx global top-$K$ & $\approx0.05$ \\
\sieve{} & per-token width in $\{0,1,2,3,4,5,6,8\}$; Eq.~\ref{eq:discrete} per KV head, 4-bit cascade score; offline per-KV-head router & 0.07--0.09 \\
\bottomrule
\end{tabular}
\end{table}

\paragraph{Deviations from the original papers.}
KIVI keeps its last partial group in full precision; we quantize it as a short
group, and the 32-token uncompressed window shared by all methods plays the
role of KIVI's residual.  KVQuant calibrates outlier thresholds and its
non-uniform codebook offline with Fisher-weighted $k$-means; we fit both online
on each prompt's own context without weighting, which can only favor it.  H2O
uses its raw accumulated score, whose causal position bias (a token is seen by
every later query) is H2O's own.  All methods use simulated quantization: the
model reads dequantized keys, and no packed kernel is involved.

\begin{table}[h]
\centering
\small
\caption{\sieve{} (router) against each baseline over the 36 valid cells.  A win or loss is a cell-mean difference above 0.05; the interval resamples prompts within each model and context.}
\label{tab:app-h2h}
% generated by figures/make_endtask_tables.py from 5 parquets
\begin{tabular}{@{}lrrrrl@{}}
\toprule
\sieve{} (router) vs. & wins & ties & losses & mean $\Delta$ & 90\% CI \\
\midrule
TurboQuant-MSE & 3 & 17 & 16 & $-0.180$ & $[-0.208, -0.152]$ \\
SnapKV & 23 & 11 & 2 & $+0.367$ & $[+0.338, +0.398]$ \\
DropKV & 25 & 8 & 3 & $+0.348$ & $[+0.314, +0.384]$ \\
Ada-KV & 21 & 9 & 6 & $+0.262$ & $[+0.236, +0.290]$ \\
LaProx & 19 & 10 & 7 & $+0.187$ & $[+0.154, +0.221]$ \\
OBCache-K + Ada-KV & 17 & 9 & 10 & $+0.163$ & $[+0.129, +0.197]$ \\
best eviction method per cell & 16 & 10 & 10 & $+0.127$ & --- \\
\bottomrule
\end{tabular}

\end{table}

\begin{table}[h]
\centering
\scriptsize
\setlength{\tabcolsep}{3pt}
\caption{Failure case: Llama-3.1-8B at 128k.  Error types are over single- and multi-key answers at $B\in\{2,3\}$ (80 per method).  \emph{Truncated}: the prediction is a proper prefix of at least three digits of the correct number.  \emph{Wrong needle}: another needle's value.}
\label{tab:app-fail}
% generated by figures/make_endtask_tables.py from 5 parquets
\begin{tabular}{@{}lrrrrrrrr@{}}
\toprule
 & \multicolumn{2}{c}{task score} & \multicolumn{3}{c}{single/multi-key answers, $B{=}2,3$} & evicted & head-output & heads to \\
\cmidrule(lr){2-3}\cmidrule(lr){4-6}
method & $B{=}2$ & $B{=}3$ & truncated & wrong needle & other & ($B{=}2$) & error ($B{=}2$) & interior \\
\midrule
TurboQuant-MSE & 0.942 & 0.982 & 0\% & 1\% & 0\% & 0\% & 0.706 & --- \\
SnapKV & 0.523 & 0.652 & 6\% & 18\% & 11\% & 75\% & 0.187 & --- \\
DropKV & 0.480 & 0.657 & 11\% & 16\% & 10\% & 75\% & 0.201 & --- \\
Ada-KV & 0.704 & 0.872 & 1\% & 15\% & 1\% & 75\% & 0.181 & --- \\
LaProx & 0.737 & 0.863 & 5\% & 12\% & 1\% & 75\% & 0.187 & --- \\
OBCache-K + Ada-KV & 0.762 & 0.902 & 4\% & 10\% & 1\% & 75\% & 0.182 & --- \\
\sieve{} interior only & 0.409 & 0.622 & 39\% & 6\% & 6\% & 65\% & 0.096 & --- \\
\sieve{} (router) & 0.442 & 0.677 & 36\% & 4\% & 9\% & 64\% & 0.094 & 99\% \\
\bottomrule
\end{tabular}

\end{table}

\begin{table}[h]
\centering
\scriptsize
\setlength{\tabcolsep}{2.4pt}
\caption{Full end-task grid: mean score over 20 prompts per cell.  Bold marks the best method in each valid row.  $\ddagger$: uncompressed score below 0.9, excluded from every aggregate.  KIVI, KVQuant and H2O are added by the paired rerun (\S\ref{app:ledger}).}
\label{tab:app-grid}
% generated by figures/make_endtask_tables.py from 5 parquets
\begin{tabular}{@{}llllrrrrrrrrr@{}}
\toprule
model & ctx & task & $B$ & FP & TurboQ. & SnapKV & DropKV & Ada-KV & LaProx & OBC+Ada & \sieve{}-int. & \sieve{} \\
\midrule
Llama-3.1-8B & 8k & multi-key & 2 & 1.00 & \textbf{1.00} & 0.20 & 0.30 & 0.25 & 0.25 & 0.40 & 0.35 & 0.30 \\
Llama-3.1-8B & 8k & multi-key & 3 & 1.00 & \textbf{1.00} & 0.35 & 0.55 & 0.55 & 0.55 & 0.75 & 0.60 & \textbf{1.00} \\
Llama-3.1-8B & 8k & multi-value & 2 & 1.00 & \textbf{1.00} & 0.14 & 0.16 & 0.19 & 0.25 & 0.39 & 0.15 & 0.23 \\
Llama-3.1-8B & 8k & multi-value & 3 & 1.00 & \textbf{1.00} & 0.34 & 0.42 & 0.42 & 0.46 & 0.59 & 0.40 & 0.99 \\
Llama-3.1-8B & 8k & single & 2 & 1.00 & \textbf{1.00} & 0.45 & 0.55 & 0.50 & 0.60 & 0.70 & 0.30 & 0.35 \\
Llama-3.1-8B & 8k & single & 3 & 1.00 & \textbf{1.00} & 0.75 & 0.65 & 0.75 & 0.80 & 0.75 & 0.60 & \textbf{1.00} \\
Llama-3.1-8B & 8k & var.\ track. & 2 & 1.00 & \textbf{0.98} & 0.45 & 0.49 & 0.42 & 0.40 & 0.54 & 0.43 & 0.60 \\
Llama-3.1-8B & 8k & var.\ track. & 3 & 1.00 & \textbf{1.00} & 0.60 & 0.64 & 0.62 & 0.59 & 0.88 & 0.67 & 0.95 \\
\addlinespace[1pt]
Llama-3.1-8B & 32k & multi-key & 2 & 1.00 & \textbf{1.00} & 0.30 & 0.20 & 0.35 & 0.45 & 0.55 & 0.40 & 0.45 \\
Llama-3.1-8B & 32k & multi-key & 3 & 1.00 & \textbf{1.00} & 0.45 & 0.35 & 0.70 & 0.70 & 0.70 & 0.60 & \textbf{1.00} \\
Llama-3.1-8B & 32k & multi-value & 2 & 1.00 & \textbf{0.97} & 0.12 & 0.06 & 0.20 & 0.26 & 0.35 & 0.16 & 0.23 \\
Llama-3.1-8B & 32k & multi-value & 3 & 1.00 & \textbf{1.00} & 0.42 & 0.23 & 0.56 & 0.61 & 0.70 & 0.40 & 0.90 \\
Llama-3.1-8B & 32k & single & 2 & 1.00 & \textbf{1.00} & 0.60 & 0.35 & 0.65 & 0.75 & 0.85 & 0.55 & 0.60 \\
Llama-3.1-8B & 32k & single & 3 & 1.00 & \textbf{1.00} & 0.70 & 0.60 & 0.95 & 0.95 & 0.90 & 0.65 & \textbf{1.00} \\
Llama-3.1-8B & 32k & var.\ track. & 2 & 1.00 & \textbf{0.99} & 0.28 & 0.37 & 0.46 & 0.47 & 0.69 & 0.46 & 0.63 \\
Llama-3.1-8B & 32k & var.\ track. & 3 & 1.00 & \textbf{1.00} & 0.65 & 0.55 & 0.82 & 0.87 & 0.86 & 0.64 & \textbf{1.00} \\
\addlinespace[1pt]
Llama-3.1-8B & 128k & multi-key & 2 & 1.00 & \textbf{1.00} & 0.30 & 0.40 & 0.50 & 0.60 & 0.65 & 0.35 & 0.40 \\
Llama-3.1-8B & 128k & multi-key & 3 & 1.00 & \textbf{0.95} & 0.45 & 0.50 & 0.80 & 0.70 & 0.85 & 0.50 & 0.55 \\
Llama-3.1-8B & 128k & multi-value & 2 & 0.97 & \textbf{0.94} & 0.41 & 0.35 & 0.57 & 0.59 & 0.64 & 0.34 & 0.34 \\
Llama-3.1-8B & 128k & multi-value & 3 & 0.97 & \textbf{0.99} & 0.54 & 0.55 & 0.79 & 0.81 & 0.85 & 0.55 & 0.69 \\
Llama-3.1-8B & 128k & single & 2 & 1.00 & \textbf{1.00} & 0.90 & 0.75 & \textbf{1.00} & 0.95 & 0.95 & 0.45 & 0.45 \\
Llama-3.1-8B & 128k & single & 3 & 1.00 & \textbf{1.00} & 0.95 & 0.85 & \textbf{1.00} & \textbf{1.00} & 0.95 & 0.65 & 0.65 \\
Llama-3.1-8B & 128k & var.\ track. & 2 & 1.00 & \textbf{0.83} & 0.48 & 0.42 & 0.74 & 0.81 & 0.81 & 0.50 & 0.58 \\
Llama-3.1-8B & 128k & var.\ track. & 3 & 1.00 & \textbf{0.99} & 0.67 & 0.73 & 0.90 & 0.94 & 0.96 & 0.79 & 0.82 \\
\addlinespace[1pt]
Qwen3-8B & 8k & multi-key & 2 & 1.00 & \textbf{0.95} & 0.00 & 0.05 & 0.00 & 0.05 & 0.05 & 0.15 & 0.90 \\
Qwen3-8B & 8k & multi-key & 3 & 1.00 & \textbf{1.00} & 0.10 & 0.20 & 0.15 & 0.30 & 0.10 & 0.50 & \textbf{1.00} \\
Qwen3-8B & 8k & multi-value$^\ddagger$ & 2 & 0.24 & 0.69 & 0.00 & 0.05 & 0.00 & 0.01 & 0.01 & 0.05 & 0.72 \\
Qwen3-8B & 8k & multi-value$^\ddagger$ & 3 & 0.24 & 0.49 & 0.00 & 0.07 & 0.00 & 0.21 & 0.10 & 0.28 & 0.59 \\
Qwen3-8B & 8k & single & 2 & 1.00 & 0.95 & 0.10 & 0.10 & 0.05 & 0.25 & 0.15 & 0.20 & \textbf{1.00} \\
Qwen3-8B & 8k & single & 3 & 1.00 & \textbf{1.00} & 0.25 & 0.35 & 0.25 & 0.55 & 0.40 & 0.60 & \textbf{1.00} \\
Qwen3-8B & 8k & var.\ track. & 2 & 1.00 & 0.46 & 0.19 & 0.28 & 0.28 & 0.41 & 0.43 & 0.47 & \textbf{0.99} \\
Qwen3-8B & 8k & var.\ track. & 3 & 1.00 & \textbf{1.00} & 0.63 & 0.63 & 0.62 & 0.66 & 0.68 & 0.81 & \textbf{1.00} \\
\addlinespace[1pt]
Qwen3-8B & 32k & multi-key & 2 & 1.00 & 0.75 & 0.00 & 0.15 & 0.05 & 0.15 & 0.15 & 0.10 & \textbf{1.00} \\
Qwen3-8B & 32k & multi-key & 3 & 1.00 & \textbf{1.00} & 0.10 & 0.30 & 0.10 & 0.30 & 0.25 & 0.40 & \textbf{1.00} \\
Qwen3-8B & 32k & multi-value$^\ddagger$ & 2 & 0.04 & 0.75 & 0.00 & 0.00 & 0.00 & 0.05 & 0.00 & 0.01 & 0.12 \\
Qwen3-8B & 32k & multi-value$^\ddagger$ & 3 & 0.04 & 0.42 & 0.00 & 0.04 & 0.00 & 0.03 & 0.00 & 0.00 & 0.07 \\
Qwen3-8B & 32k & single & 2 & 1.00 & 0.95 & 0.20 & 0.30 & 0.20 & 0.40 & 0.25 & 0.35 & \textbf{1.00} \\
Qwen3-8B & 32k & single & 3 & 1.00 & \textbf{1.00} & 0.30 & 0.55 & 0.50 & 0.70 & 0.55 & 0.45 & \textbf{1.00} \\
Qwen3-8B & 32k & var.\ track. & 2 & 1.00 & 0.40 & 0.28 & 0.40 & 0.45 & 0.79 & 0.61 & 0.46 & \textbf{0.99} \\
Qwen3-8B & 32k & var.\ track. & 3 & 1.00 & 0.95 & 0.69 & 0.70 & 0.77 & 0.92 & 0.81 & 0.76 & \textbf{0.99} \\
\bottomrule
\end{tabular}

\end{table}

\section{Estimate ledger}
\label{app:ledger}

Numbers typeset as \est{x} in the main text are estimates, not measurements.
Each is replaced by the run named below; the paired rerun (R12-A) also
replaces every measured end-task number, because the forward pass is not
bit-reproducible across processes and all methods must share one process per
cell to be compared prompt by prompt.

\begin{table}[h]
\centering
\scriptsize
\setlength{\tabcolsep}{3pt}
\begin{tabular}{@{}p{0.2\linewidth}p{0.13\linewidth}p{0.49\linewidth}p{0.1\linewidth}@{}}
\toprule
quantity & estimate & basis & replaced by \\
\midrule
H2O, Llama / Qwen & 0.43 / 0.22 & Question-aware H2O at Llama-3.1-8B 32k (two runs): 0.37 at $B=2$, 0.48 at $B=3$.  H2O's score sums attention from every prefill query, so the 32 question rows barely change it; Qwen scaled by the SnapKV Llama/Qwen ratio. & R12-A \\
KIVI ($g{=}128$), Llama / Qwen & 0.98 / 0.95 & Llama: TurboQuant is at 0.985 and KIVI has lower head-output error at equal code bits on Llama-3.2-1B (0.45 vs.\ 0.71 at $B=2$).  Qwen: on Qwen3-1.7B at 2k (11 prompt-task rows, prompts 2000--2003) TurboQuant scores 0/11 at both $B=2$ and $B=3$; KIVI-$g128$ scores 0.95 at $B=2$ and 1.00 at $B=3$. & R12-A \\
KVQuant, Llama / Qwen & 0.99 / 0.98 & As KIVI; KVQuant scores 11/11 at both budgets on the Qwen3-1.7B check and has the lowest head-output error on Llama-3.2-1B (0.21 at $B=2$). & R12-A \\
Hard cell, TurboQuant & 0.83 & Measured on a screening block (prompts 500--539, 40 prompts): FP 1.00, TurboQuant 0.825, $\Delta$ 90\% CI $[0.075,0.275]$.  The evaluation uses fresh prompts 1000--1059. & R12-B \\
Hard cell, other methods & see Table~\ref{tab:endtask} & The same methods on 4-key retrieval at 32k, $B=2$, scaled by 0.7 for eight times as many needles.  KIVI and KVQuant: TurboQuant's screening value plus a small margin for their lower output error. & R12-B \\
Pooled-score interior, 128k, $B=2$ & 0.58 & Just below midway between the interior (0.409) and OBCache-K + Ada-KV (0.762); the pre-registered ``supported'' threshold is 0.586.  A placeholder, not a prediction of the test's outcome. & R12-A (diag.) \\
Oracle router, 128k, $B=2$ & 0.55 & At 32k, $B=2$ the oracle adds $+0.30$ over the calibrated router; scaled to 128k, where the calibrated router already sends 99\% of heads to the interior. & R12-A (diag.) \\
\bottomrule
\end{tabular}
\end{table}
